\section{Introduction}\label{sec:introduction}
\suppressfloats[t]

The Birkhoff conjecture asks whether an integrable planar convex billiard
must be an ellipse. We prove the version in which integrability is assumed
only in a neighborhood of the boundary, with continuous dependence of the
caustics. The individual leaves and the billiard boundary are smooth;
no differentiability of the foliation across its leaves is required.

\begin{definition}[A continuous caustic collar]\label{def:caustic-collar}
Let \(\Omega\subset\R^2\) be a bounded strictly convex domain whose boundary
is a \(C^\infty\) simple closed regular curve with positive curvature.
A \emph{smooth closed convex caustic} is a \(C^\infty\) simple closed
regular curve in \(\Omega\), bounding a strictly convex set, such that a
billiard segment tangent to the curve remains tangent after reflection
at either endpoint. Here reflection of an incoming unit velocity \(v\)
at a boundary point with outward unit normal \(n\) is
\[
v\longmapsto v-2(v\cdot n)n.
\]
A \emph{continuous caustic collar} consists of a relatively open set
\(A\subset\overline\Omega\) containing \(\partial\Omega\) and a homeomorphism
\[
H_0:\mathbb S^1\times[0,1)\longrightarrow A
\]
such that \(H_0(\mathbb S^1\times\{0\})=\partial\Omega\), and
\(H_0(\mathbb S^1\times\{b\})\) is a smooth closed convex caustic for every
\(0<b<1\).
\end{definition}

\begin{theorem}\label{thm:main}
If \(\Omega\) has a continuous caustic collar, then \(\Omega\) is an ellipse.
Equivalently, there are \(c\in\R^2\) and a real symmetric positive-definite
matrix \(Q\) such that
\[
\Omega=\{x\in\R^2:(x-c)^TQ(x-c)<1\}.
\]
Circles are included.
\end{theorem}

This resolves the near-boundary Birkhoff conjecture positively in the
smooth positively curved class of \Cref{def:caustic-collar}. The collar
contains every intermediate leaf. A Cantor family of invariant curves,
even one accumulating at the boundary, does not satisfy this assumption.
The hypothesis makes no assertion about the billiard dynamics farther
inside the domain.

\begin{figure}[t]
\centering
\begin{minipage}{0.48\textwidth}
\centering
\begin{tikzpicture}[scale=.85,>=Stealth]
\fill[deepblue!9,even odd rule] (0,0) ellipse (3 and 2)
 (0,0) ellipse (2.70 and 1.5133);
\draw[thick] (0,0) ellipse (3 and 2);
\draw[deepblue] (0,0) ellipse (2.90 and 1.8466);
\draw[deepblue] (0,0) ellipse (2.80 and 1.6852);
\draw[deepblue] (0,0) ellipse (2.70 and 1.5133);
\fill (-2.236,0) circle (1.4pt);
\fill (2.236,0) circle (1.4pt);
\node[font=\small] at (0,-2.38) {A full caustic collar};
\node[font=\small,anchor=west] at (2.35,1.6) {$\partial\Omega$};
\end{tikzpicture}
\end{minipage}\hfill
\begin{minipage}{0.48\textwidth}
\centering
\begin{tikzpicture}[scale=1.42,>=Stealth]
\draw[thick,domain=-1.3:1.3,samples=40,variable=\y]
 plot ({1-.3*\y*\y},{\y});
\draw[dashed,gray] (1,-1.28)--(1,1.28);
\draw[->,deepblue,thick] (.47,-1.12)--(1,0);
\draw[->,deepblue,thick] (1,0)--(.47,1.12);
\draw[->] (1,0)--(2.05,0) node[right,font=\small] {$n_\psi$};
\draw[->,gray] (1,0)--(1.95,-.44);
\draw[->,gray] (1,0)--(1.95,.44);
\draw (1.56,0) arc (0:25:.56);
\draw (1.70,-.326) arc (-25:0:.77);
\node[font=\small] at (1.80,.24) {$d$};
\node[font=\small] at (1.93,-.17) {$d$};
\node[font=\small,anchor=east] at (.43,-.88) {$e_{\psi-d}$};
\node[font=\small,anchor=east] at (.43,.88) {$e_{\psi+d}$};
\fill (1,0) circle (1.3pt);
\node[font=\small,anchor=east] at (.94,.14) {$\Gamma$};
\node[font=\small] at (1,-1.60) {Reflection in line coordinates};
\end{tikzpicture}
\end{minipage}
\caption{The geometric hypothesis and the local reflection convention.
The left panel shows confocal elliptic caustics, the model covered by
the theorem. The shaded collar represents a full family of leaves.
At the impact point \(\Gamma\) in the right panel, \(n_\psi\) is the
outward unit normal. The blue arrows give the incoming and outgoing
unit-velocity directions \(e_{\psi-d}\) and \(e_{\psi+d}\); their line
normals have angles \(\psi-d\) and \(\psi+d\).}
\label{fig:caustic}
\end{figure}


\subsection{History and scope}

Ellipses are the basic integrable examples: their confocal conics supply
caustics, and a confocal elliptic family fills a collar of the boundary.
The conjecture is traditionally attributed to Birkhoff; its early
printed formulation appears in Poritsky's work~\cite{Poritsky1950}.
For the classical dynamical setting and this historical attribution,
see \cite{Birkhoff1927,KS2018}. Poritsky's translation property for
nested string constructions is stronger than the existence of a
foliation~\cite{Poritsky1950}. Lazutkin constructed a positive-measure
family of grazing caustics for sufficiently smooth positively curved
tables~\cite{Lazutkin1973}; its gaps do not give a full collar.
Marvizi and Melrose's interpolating grazing integral is accurate up to
a flat error, rather than an exact integral~\cite{MM1982}. Mather's
obstruction when curvature vanishes underscores the positive-curvature
hypothesis here~\cite{Mather1982}.

Several rigidity theorems establish parts of the conjecture under
additional assumptions. Bialy's rigidity theorem assumes a foliation
of the whole phase cylinder by continuous noncontractible invariant
curves and concludes that the table is a
circle~\cite{Bialy1993}. Innami's theorem concerns convex caustics
whose rotation numbers approach \(1/2\), the opposite endpoint from
grazing rotation zero; Arnold and Bialy give a geometric treatment of
that result~\cite{Innami2002,AB2018}. Bialy and Mironov prove ellipse
rigidity for centrally symmetric tables when the phase annulus reaches
the four-periodic curve, with a corresponding \(C^1\) foliation by
physical caustics through rotation \(1/4\)~\cite{BM2022}. These
hypotheses are not supplied by a small grazing collar.

Perturbative results include the local rigidity work of Avila, De
Simoi and Kaloshin near low-eccentricity ellipses~\cite{ADK2016}, and
Kaloshin and Sorrentino near ellipses of arbitrary
eccentricity~\cite{KS2018}. Kaloshin, Koudjinan and Zhang treat nearly
centrally symmetric domains under specified caustic
conditions~\cite{KKZ2024}; Koval treats small rational-rotation
caustics near almost every ellipse~\cite{Koval2026}. Glutsyuk's
classification for complex algebraic caustics lies in a different
analytic setting~\cite{Glutsyuk2026}.

\Cref{thm:main} assumes none of those perturbative, symmetry, or algebraic
conditions. In particular, we must obtain analytic regularity from the
given smooth geometry and merely continuous foliation. The proof below
does so before using complex continuation. Its grazing normalization
and action expansions have antecedents in Lazutkin and Marvizi--Melrose;
its rational-action Fourier test is related to Avila--De Simoi--Kaloshin;
and its use of complex singularities has a perturbative predecessor in
Kaloshin--Sorrentino~\cite{Lazutkin1973,MM1982,ADK2016,KS2018}.
These are method ancestry, not formal inputs to the arbitrary-table
analytic inversion, first-wall, or probability arguments proved here.

\subsection{The main ideas}

The proof has three parts. First, rational caustics force analytic
regularity. Nesting and area recurrence imply that every sufficiently
small rational rotation occurs on a curve entirely made of periodic
orbits. The actions of these orbits give a high-frequency system whose
derivative is the identity up to a small operator. An analytic Newton
construction then upgrades the boundary from smooth to analytic.
A second construction uses both values and first parameter jets at
rational resonances to obtain a jointly analytic conjugacy at glancing.

Second, finite reflected chains determine their intermediate lines
algebraically from endpoint jets. The proof combines a finite system
for one-variable jets with a coalescing-impact scaling. This local
determination result controls monodromy: a finite first continuation
wall cannot have even one local joint meromorphic patch. The argument
separates finite jet determination from global continuation, and the
corresponding lemmas are stated independently of the later growth
analysis.

Third, we study the complexified boundary through its meromorphic
normal and two analytic derivative fields. Zeros or poles of that
normal force an elliptic form. If there are none, continued impacts
admit one complex shape coordinate. Growth estimates exclude most
possible behaviors at the edge of its strip, leaving a critical
finite-width case. There, moment bounds and paths through regions
where the fields are large limit the total angular change, an
\emph{angle budget}. For one fixed odd periodic chain, matrix identities
transfer that control to Brownian paths stopped before the shape
boundary. Finally, we find a horizontal row on which the first field is
robustly small or large outside a set of vanishing relative measure.
The sign test on this row keeps the last moving impact away from the
boundary, contradicting the assumed first wall.

Two aspects of these arguments may be useful separately. The first is
the finite-chain algebraic determination principle and its monodromy
consequence. The second is the combination of an individual-center
quiet-profile obstruction with moment bounds: a local analytic passage
forces a global width improvement, so even an exceptional sequence of
quiet centers is ruled out. The subsequent probability arguments use
this pointwise conclusion explicitly.

\Needspace{20\baselineskip}
\subsection{Organization and dependencies}

The sections and their main deductions are recorded below. The
arbitrary-input action estimates of \Cref{ana:section} are separate from
the geometric vanishing supplied by the caustic collar. The parameter
jets in \Cref{joint:section} use the full collar again.

\begin{center}
\begin{tabular}{@{}p{0.24\textwidth}p{0.69\textwidth}@{}}
\toprule
Location & Role in the proof \\
\midrule
\Cref{ana:section,joint:section}
& Rational-action inversion gives an analytic boundary; rational jets give a jointly analytic glancing parametrization. \\
\Cref{sec:continuation}
& Finite-chain algebraic determination rules out a local patch at a first continuation wall. \\
\Cref{sec:shapes}
& Zeros or poles of the normal force an ellipse; otherwise impacts lift to a global shape coordinate. \\
\Cref{sec:cones,sec:profiles}
& Kernel continuation and profile obstructions reduce growth to a critical finite wall. \\
\Cref{sec:critical,sec:model}
& Conditional critical estimates precede the independent exclusion of prohibited sequences; only then is the angle budget activated. \\
\Cref{sec:odd-slice,sign:section}
& Stopped odd-chain geometry and a sign row exclude the critical wall and complete \Cref{thm:main}. \\
\bottomrule
\end{tabular}
\end{center}

There is one deliberate forward reference: the angle-budget reduction
uses the exclusion of a precisely stated class of sequences, proved in
the following section. That exclusion uses the earlier moment and
continuation estimates, not the angle budget. The final proof of
\Cref{thm:main} assembles the alternatives after \Cref{thm:last-site}.

\subsection{Conventions}

Angles are real lifts of \(2\pi\)-periodic variables. Complex Euclidean
products are bilinear, without complex conjugation. A meromorphic
projective map means meromorphic affine coordinates in a generic chart;
all incidence identities are used on their regular locus and then
continued. Constants denoted by \(C\) may increase between estimates.
Dependence on a fixed aperture or another fixed auxiliary parameter is
displayed when it matters for an order of choices.

The normalized generating action is denoted by \(D\). Later, an uppercase
\(D\) also denotes a positive rescaled depth inside a local scaling
argument; its meaning is stated there. To keep the width bookkeeping
distinct, the phase-width deficit is denoted throughout by
\(\Delta(t)=H-h_X(t)\). Horizontal expectations are normalized averages
over one period unless a different probability law is specified.
